\documentclass[12pt,a4paper]{exam}

% --------------------
% Packages
% --------------------
\usepackage[utf8]{inputenc}
\usepackage[T1]{fontenc}
\usepackage[french]{babel}
\usepackage{amsmath,amssymb}
\usepackage{geometry}
\usepackage{pgfplots}
\pgfplotsset{compat=1.18}
\geometry{margin=2cm}

% --------------------
% En-tête / pied (exam)
% --------------------
\header{}{DS de mathématiques}{1ère STMG}
\footer{}{\thepage}{}

% --------------------
\begin{document}

\begin{center}
    \Large\textbf{CORRECTION -- Interrogation\\Dérivation}
\end{center}

\vspace{0.3cm}

\noindent
\textbf{Durée :} 1h \hfill
\textbf{Calculatrice :} \underline{NON} autorisée

\vspace{0.5cm}

% ========================================================
%             Remarques
% ========================================================
{\color{red}
\underline{\textbf{Remarques générales :}}

\vspace{5pt}
\begin{itemize}
    \item les éléments donnés dans cette correction sont simplement les méthodes à mettre 
    en place pour répondre aux questions. À vous de refaire les calculs et la démarche associée
    \item \underline{Exercice 1 :}
    \begin{itemize}
        \item beaucoup de confusions entre $f'(0)$, $f'(3)$ et $f(0)$, $f(3)$. Là on veut f' donc 
        cela correspond à la pente des tangente aux points donnés. Pour $f'(0)$ c'est la 
        pente de la tangente au point $x=0$, donc de la tangente en $A$. Il ne s'agit pas du calcul
        de $f(0)$ qui lui est une simple lecture d'image : $f(0) = 2$.
        \item ici beaucoup m'ont mis le calcul de la pente pour $T_A$ et $T_B$. Or, il s'agit de trouver 
        l'équation de la tangente qui est d'après le cours : $y = f'(a)x + p$ et il faut alors déterminer $f'(a)$
        (fait à la question d'avant) et $p$.
    \end{itemize}
    \item \underline{Exercice 2 :}
    \begin{itemize}
        \item il faut absolument retravailler les tableaux de signes des fonctions affines, 
        réussis par très peu d'entre vous
        \item le tableau de variations découle directement de ce tableau de signes
    \end{itemize}
\end{itemize}
}

\vspace{10pt}
% ======================================================
\noindent
\textbf{Exercice 1 : Dérivées et tangentes (\ \ \ \  / 5 points )}
\vspace{10pt}
\noindent\\On considère la fonction $f$ représentée ci-dessous ainsi que deux tangentes à $f$.

\vspace{-5pt}
\begin{center}
    \begin{tikzpicture}[scale=1, xscale=1.5]

    % Grille
    \draw[step=1cm, gray, very thin] (-3,-2) grid (6,7);

    % Axes
    \draw[->] (-3.2,0) -- (6.2,0) node[right] {$x$};
    \draw[->] (0,-2.2) -- (0,7.2) node[above] {$y$};

    % Graduations x
    \foreach \x in {-3,-2,-1,1,2,3,4,5,6}
        \draw (\x,0.1) -- (\x,-0.1) node[below] {\small $\x$};

    % Graduations y
    \foreach \y in {-1,1,2,3,4,5,6}
        \draw (0.1,\y) -- (-0.1,\y) node[left] {\small $\y$};

    % Courbe y = 3x^2 - 4 (limitée à y ≤ 5)
    \draw[domain=-1.5:3.5, samples=200, smooth, thick, black]
        plot (\x, {\x*\x - 2*\x + 2});

    \node at (-1,6.2) {$\mathcal{C}_f$};

    % Courbe y = 3x^2 - 4 (limitée à y ≤ 5)
    \draw[domain=-2.5:2, samples=200, smooth, very thick, black]
        plot (\x, {-2*\x + 2});

    \node at (-1.5,4.5) {$\mathcal{T}_A$};

    % Courbe y = 3x^2 - 4 (limitée à y ≤ 5)
    \draw[domain=1.25:3.5, samples=200, smooth, very thick, black]
        plot (\x, {4*\x -7});

    \node at (2.5,1.75) {$\mathcal{T}_B$};

    \filldraw (0, 2) circle (2pt) node[above right] {A};

    \filldraw (3, 5) circle (2pt) node[below right] {B};

    \end{tikzpicture}
\end{center}


\begin{questions}

\question[2] Déterminer graphiquement $f'(0)$ et $f'(3)$.

{\color{red}Ici, il s'agit de $f'$ donc on calcule les pentes des tangentes aux point donnés : $x=0$ et $x=3$

On trouve : $f'(0)=-2$ et $f'(3) = 4$.

Je compte sur vous pour refaire le calculs si vous n'avez pas les bonnes réponses (avec l'aide du cours si besoin).}

\question[3] En déduire les équations des tangentes $\mathcal{T}_A$ et $\mathcal{T}_B$ à 
la courbe $\mathcal{C}_f$.

{\color{red}D'après le cours une équation de tangente est : $y = f'(a)x + p$

Nous avons déjà calculé $f'(a)$ pour nos deux tangentes : $f'(0)=-2$ et $f'(3) = 4$.

On a donc $T_A : y = -2x+p$ et $T_B : y = 4x + p$.

Il nous reste plus qu'à calculer $p$ : c'est l'ordonnée à l'origine (évident pour $T_A$ qui vaut $2$). Pour 
$T_B$ il faut utiliser la méthode du point qui appartient à la droite comme vu dans le cours.

On trouve $T_A : y = -x+2$ et $T_B : y = 3x - 7$.
}

\end{questions}

\vspace{1cm}

% ========================================================
%             Exercice 2
% ========================================================
\noindent
\textbf{Exercice 2 : Calculs de dérivées ( \ \ \ / 5 points )}

\vspace{10pt}
\noindent

\begin{questions}

\question[2] Déterminer l'expression de la fonction dérivée de chacune des fonctions suivantes :

\vspace{5pt}
\hspace{1cm}a) $f(x) = 3x^3 - 7x$ \hspace{1cm}b) $g(x) = 3x^2 - x + 5$

{\color{red}a) $f'(x) = 9x^2 - 7$

b) $f'(x) = 6x-1$}

\vspace{10pt}

\question On considère la fonction : $f(x) = 2x^2 + 4x - 2$.

\vspace{5pt}
\hspace{1cm}a) (1 point) Calculer la dérivée $f'(x)$.

{\color{red}$f'(x)=4x + 4$}

\vspace{5pt}
\hspace{1cm}b) (1 point) Étudier le signe de cette dérivée.

{\color{red}
Il faut revoir cette question pour la majorité de la classe.

Je rappelle la méthode : on cherche la valeur ou ça change de signe en résolvant l'équation : 
$4x + 4 = 0$.

Ensuite on construit le tableau de signes.

Pour remplire les bons signes aux bons endroits il suffit de regarder le signe du coefficient directeur : ici $4$.

Comme $4$ est positif alors ce sera d'abord un signe $-$ puis un signe $+$ dans le tableau. Sinon c'est l'inverse.

REFAITES CET EXERCICE.
}

\vspace{5pt}
\hspace{1cm}c) (1 point) Dresser le tableau de variation de la fonction $f$.

{\color{red}
Le tableau de variations découle directement du tableau de signe. Revoir le cours de dérivées si besoin.
}

\end{questions}

\vspace{0.5cm}

\begin{center}
    \textit{Fin du devoir}
\end{center}

\end{document}
